% Generated draft: requires mathematical and source review.
\documentclass[11pt]{article}
\usepackage[margin=25mm]{geometry}
\usepackage{fontspec}
\setmainfont{Latin Modern Roman}
\usepackage{amsmath,amssymb,mathtools,unicode-math}
\setmathfont{Latin Modern Math}
\usepackage{enumitem,needspace,xurl,hyperref}
\setlength{\parindent}{0pt}
\setlength{\parskip}{5pt}
\setlength{\emergencystretch}{3em}
\newcommand{\N}{\mathbb{N}}
\newcommand{\Z}{\mathbb{Z}}
\newcommand{\Q}{\mathbb{Q}}
\newcommand{\R}{\mathbb{R}}
\newcommand{\C}{\mathbb{C}}
\newcommand{\F}{\mathbb{F}}
\newcommand{\A}{\mathbb{A}}
\newcommand{\PP}{\mathbb{P}}
\newcommand{\E}{\mathbb{E}}
\newcommand{\eps}{\varepsilon}
\newcommand{\dd}{\,\mathrm{d}}
\begin{document}
\section*{UnsolvedMath ID: 30004430}
\textbf{Source problem code:} catalogue entry
\subsection{Definitions and mathematical statement}
\textbf{Supplied question.} What conditions on the size or shape of the state set \(D\) ensure the existence of a \(g\)-measure?

\textbf{Unresolved definitions and quantifiers.} The supplied record calls \(D\) a state set but specifies neither its ambient space nor its measurable, topological, or geometric structure. It does not define \(g\), its domain, its normalization or regularity assumptions, or the meaning of a \(g\)-measure. Consequently, neither ``size'' nor ``shape'' has a specified mathematical meaning here. It is also unspecified whether \(g\) is fixed, ranges over a prescribed class, or is to be chosen together with the measure. These omissions prevent a precise quantified existence statement. The faithful target is therefore the classification question above, with the definitions, hypotheses, and quantifiers still requiring recovery from the underlying source; no particular interpretation of \(D\) or convention for \(g\)-measures is imposed.
\subsection{Short English statement}
Determine which size or shape conditions on \(D\) guarantee a \(g\)-measure. The supplied extract leaves the state space, the function \(g\), and the intended compatibility conditions undefined, so a precise formulation requires further source clarification. The supplied curation labels this an open classification problem.
\subsection{Sources}
\begin{itemize}
\item Supplied local UnsolvedMath registry record, permanent ID \(30004430\), source problem code \texttt{catalogue entry}, titled ``Existence Criteria for \(g\)-Measures,'' in the Oberwolfach Reports Open Problems collection. No individual proposer is identified.
\item Supplied report citation: \emph{Mini-Workshop: One-sided and Two-sided Stochastic Descriptions} (2021), Oberwolfach Reports, pp.\ 601--637, DOI \texttt{10.4171/owr/2020/11}. \url{https://doi.org/10.4171/owr/2020/11}. The record gives proposed year 2020; the citation's year 2021 and the 2020 identifiers are retained without reconciliation.
\item Supplied literature source, labeled OWR-2020-11: \url{https://ems.press/content/serial-article-files/46847}.
\item Literature assessment dated \texttt{2026-08-22}: the supplied curation describes the question as an open classification problem asking for general size and shape conditions on \(D\), rather than a single resolved theorem. This is supplied dated curation, not an independently verified literature assessment. The record also labels the statement \texttt{corrected\_verified} and reports a manual semantic, typographic, and LaTeX review; those are supplied review claims and do not resolve the missing mathematical definitions.
\end{itemize}
\end{document}
